\subsubsection{数据预处理}

问题三在问题二的历史实测数据之外新增附件3的光伏功率预报。两类数据记录口径不同，我们先统一时间和功率单位，再进行融合。

第一步是预报数据的时间细化。附件3按0:00、6:00、12:00、18:00四个发布时刻分别给出未来24小时的逐小时光伏功率预报，而模型的时间步长为10分钟，故须作尺度细化。我们以发布时刻最近已完成时段的实测功率$v_0$为起点，将逐小时预报值作为节点作线性插值，得到10分钟粒度的预报曲线：
\begin{equation}\label{eq:src30}
    Y_{s,t}=\left[Y^{(h)}+\bigl(Y^{(h+1)}-Y^{(h)}\bigr)\theta\right]\Delta t,
 \quad\theta=\frac{u_t-h}{1\text{ h}},\quad h\le u_t<h+1,
\end{equation}
其中，$Y^{(h)}$为发行时刻之后第$h$小时的预报功率（kW），$Y^{(0)}=v_0$；$u_t$为目标时段结束时刻相对发行时刻的小时数，$Y_{s,t}$为插值后乘以$\Delta t$的时段电量（kWh）。未被该次预报覆盖的时段由历史预测补齐。

第二步是预报偏差的标定。将细化后的预报曲线与附件2的光伏实测序列逐时段作差，得到以发布时刻与小时编号为索引的偏差样本：
\begin{equation}\label{eq:src31}
    z_{s,t}^{(d)}=V_t^{(d)}-Y_{s,t}^{(d)},\qquad s\in\{0,1,2,3\},\quad t\ge36s
\end{equation}
其中，$s$为发布序号，对应时刻为$6s$时；$z_{s,t}^{(d)}$为第$d$天该次发布对第$t$时段的预报偏差（kWh）。我们按小时$h=\lfloor t/6\rfloor$汇总六个时段的误差样本。仅在有效时段已经完成后登记误差；官方产品的校正误差与最终融合净负荷误差分别保存，后者用于滚动筛选。

第三步是融合预报的构造。记历史统计预报为$\hat{V}_t^{\text{hist}}$（由式\eqref{eq:src14}的同时段基准加偏差修正给出），官方预报为$Y_t$，二者按权重$\omega_t$融合：
\begin{equation}\label{eq:src32}
    \hat V_t^{\mathrm{fus}}=\max\{0,\omega_t(Y_{s,t}+b^O_{s,h})+(1-\omega_t)\hat V_t^{\mathrm{hist}}\},\qquad
 \omega_t=\frac{\mathrm{MAE}_{H,s,h}}{\mathrm{MAE}_{O,s,h}+\mathrm{MAE}_{H,s,h}+10^{-9}}.
\end{equation}
其中，$b^O_{s,h}$为对应官方预报历史误差的小时中位数，两项MAE分别评价官方与历史预测。样本不足时日间权重取0.6；0--6时及19--24时权重取0.5、不作官方小时偏差校正。日间权重随已实现误差动态更新，各时段均保留原始光伏观测。

第四步是分位数带的构造。以同发布序号、同小时内六个时段的净负荷历史误差样本为基础，给出0.70与0.90两个分位数所对应的上下边界：
\begin{equation}\label{eq:src33}
    B_{t}^{\text{lo}}=\hat{N}_{t}+e_{t,0.70},\qquad B_{t}^{\text{hi}}=\hat{N}_{t}+e_{t,0.90}
\end{equation}
该分位数带用于第四种备选合同的方向约束。模型在6、12、18时根据当前已发布预报和已完成观测复核合同。

经上述处理，我们得到两类输入：日前滚动规划所需的预测序列，以及用于筛选合同的历史误差路径。



\paragraph{组合预测的估计机制与更新顺序}
历史光伏预测包含七日加权同时段均值与近28日中位数偏差两部分。前者保留日内出力形态并增加近期观测的权重，后者修正持续高估或低估。中位数是绝对偏差损失的最小化解：
\begin{equation}\label{eq:median_fit}
 b^* = \arg\min_b\sum_{e\in\mathcal R}|e-b|.
\end{equation}
官方预报的小时偏差采用相同估计准则。式\eqref{eq:src32}的组合权重使历史误差较小的预测源获得较大权重；每次权重更新均读取已经完成时段的误差档案。

对第$d$天，我们先用此前观测计算历史基准及校正量；在0时融合新发布预报并冻结日前预测，生成原计划；在6、12、18时更新未执行窗口并生成备选合同；当日观测完成后，将预测误差归入误差档案，跨日误差在次日相应时段完成后登记。各发布时点均先更新信息、再提交合同；实时执行只使用截至当前已实现的观测，满足滚动预测控制的时序约束\cite{parisio2014mpc}。
